\subsection{Sprint 4}

Na \textit{sprint} final, comecei resolvendo o problema crítico de memória no servidor do \textit{backend}, identificado na \textit{sprint} passada e que estava impactando diretamente os processos de \textit{build} e execução. Em seguida, finalmente registrei o GitLab \textit{Runner} diretamente nesse servidor para consolidar o fluxo de \textit{deploy} do \textit{backend}.
 
A partir daí, iniciei uma nova série de depurações relacionadas ao \textit{Runner}: surgiram múltiplos problemas de permissão para a execução de comandos no nosso servidor, além de erros envolvendo operações do Git e falhas ao executar \textit{git fetch}. O processo ainda esbarrou em um erro final de falta de memória, impedindo o \textit{build} de prosseguir até que os ajustes necessários fossem feitos.
 
Após superar essas barreiras, o \textit{deploy} automático do \textit{backend} via \textit{Runner} finalmente ficou pronto e funcional. Nesse período, continuei tentando automatizar o \textit{deploy} do \textit{frontend} utilizando \textit{webhooks}, mas, apesar das tentativas e do tempo que tínhamos, essa etapa não foi concluída no projeto.
\\

\fbox{%
    \parbox{0.95\linewidth}{%
        \setlength{\parindent}{15pt}
            \itshape Por fim, depois de longas sessões de depuração, posso dizer que não entendo por que temos uma instância própria do GitLab hospedada. Caso não utilizássemos uma instância própria, teríamos muito menos dificuldades.
            Justamente por causa dessas longas sessões, somente conseguimos automatizar o \textit{deploy} do \textit{backend}, diferentemente do \textit{frontend}.
    }%
}
